% Options for packages loaded elsewhere
\PassOptionsToPackage{unicode}{hyperref}
\PassOptionsToPackage{hyphens}{url}
%
\documentclass[
  10pt,
]{article}
\usepackage{amsmath,amssymb}
\usepackage{iftex}
\ifPDFTeX
  \usepackage[T1]{fontenc}
  \usepackage[utf8]{inputenc}
  \usepackage{textcomp} % provide euro and other symbols
\else % if luatex or xetex
  \usepackage{unicode-math} % this also loads fontspec
  \defaultfontfeatures{Scale=MatchLowercase}
  \defaultfontfeatures[\rmfamily]{Ligatures=TeX,Scale=1}
\fi
\usepackage{lmodern}
\ifPDFTeX\else
  % xetex/luatex font selection
  \setmainfont[]{Noto Sans CJK JP}
  \setmonofont[]{Noto Sans Mono CJK JP}
\fi
% Use upquote if available, for straight quotes in verbatim environments
\IfFileExists{upquote.sty}{\usepackage{upquote}}{}
\IfFileExists{microtype.sty}{% use microtype if available
  \usepackage[]{microtype}
  \UseMicrotypeSet[protrusion]{basicmath} % disable protrusion for tt fonts
}{}
\makeatletter
\@ifundefined{KOMAClassName}{% if non-KOMA class
  \IfFileExists{parskip.sty}{%
    \usepackage{parskip}
  }{% else
    \setlength{\parindent}{0pt}
    \setlength{\parskip}{6pt plus 2pt minus 1pt}}
}{% if KOMA class
  \KOMAoptions{parskip=half}}
\makeatother
\usepackage{xcolor}
\usepackage[margin=0.75in]{geometry}
\usepackage{color}
\usepackage{fancyvrb}
\newcommand{\VerbBar}{|}
\newcommand{\VERB}{\Verb[commandchars=\\\{\}]}
\DefineVerbatimEnvironment{Highlighting}{Verbatim}{commandchars=\\\{\}}
% Add ',fontsize=\small' for more characters per line
\newenvironment{Shaded}{}{}
\newcommand{\AlertTok}[1]{\textcolor[rgb]{1.00,0.00,0.00}{\textbf{#1}}}
\newcommand{\AnnotationTok}[1]{\textcolor[rgb]{0.38,0.63,0.69}{\textbf{\textit{#1}}}}
\newcommand{\AttributeTok}[1]{\textcolor[rgb]{0.49,0.56,0.16}{#1}}
\newcommand{\BaseNTok}[1]{\textcolor[rgb]{0.25,0.63,0.44}{#1}}
\newcommand{\BuiltInTok}[1]{\textcolor[rgb]{0.00,0.50,0.00}{#1}}
\newcommand{\CharTok}[1]{\textcolor[rgb]{0.25,0.44,0.63}{#1}}
\newcommand{\CommentTok}[1]{\textcolor[rgb]{0.38,0.63,0.69}{\textit{#1}}}
\newcommand{\CommentVarTok}[1]{\textcolor[rgb]{0.38,0.63,0.69}{\textbf{\textit{#1}}}}
\newcommand{\ConstantTok}[1]{\textcolor[rgb]{0.53,0.00,0.00}{#1}}
\newcommand{\ControlFlowTok}[1]{\textcolor[rgb]{0.00,0.44,0.13}{\textbf{#1}}}
\newcommand{\DataTypeTok}[1]{\textcolor[rgb]{0.56,0.13,0.00}{#1}}
\newcommand{\DecValTok}[1]{\textcolor[rgb]{0.25,0.63,0.44}{#1}}
\newcommand{\DocumentationTok}[1]{\textcolor[rgb]{0.73,0.13,0.13}{\textit{#1}}}
\newcommand{\ErrorTok}[1]{\textcolor[rgb]{1.00,0.00,0.00}{\textbf{#1}}}
\newcommand{\ExtensionTok}[1]{#1}
\newcommand{\FloatTok}[1]{\textcolor[rgb]{0.25,0.63,0.44}{#1}}
\newcommand{\FunctionTok}[1]{\textcolor[rgb]{0.02,0.16,0.49}{#1}}
\newcommand{\ImportTok}[1]{\textcolor[rgb]{0.00,0.50,0.00}{\textbf{#1}}}
\newcommand{\InformationTok}[1]{\textcolor[rgb]{0.38,0.63,0.69}{\textbf{\textit{#1}}}}
\newcommand{\KeywordTok}[1]{\textcolor[rgb]{0.00,0.44,0.13}{\textbf{#1}}}
\newcommand{\NormalTok}[1]{#1}
\newcommand{\OperatorTok}[1]{\textcolor[rgb]{0.40,0.40,0.40}{#1}}
\newcommand{\OtherTok}[1]{\textcolor[rgb]{0.00,0.44,0.13}{#1}}
\newcommand{\PreprocessorTok}[1]{\textcolor[rgb]{0.74,0.48,0.00}{#1}}
\newcommand{\RegionMarkerTok}[1]{#1}
\newcommand{\SpecialCharTok}[1]{\textcolor[rgb]{0.25,0.44,0.63}{#1}}
\newcommand{\SpecialStringTok}[1]{\textcolor[rgb]{0.73,0.40,0.53}{#1}}
\newcommand{\StringTok}[1]{\textcolor[rgb]{0.25,0.44,0.63}{#1}}
\newcommand{\VariableTok}[1]{\textcolor[rgb]{0.10,0.09,0.49}{#1}}
\newcommand{\VerbatimStringTok}[1]{\textcolor[rgb]{0.25,0.44,0.63}{#1}}
\newcommand{\WarningTok}[1]{\textcolor[rgb]{0.38,0.63,0.69}{\textbf{\textit{#1}}}}
\usepackage{longtable,booktabs,array}
\usepackage{calc} % for calculating minipage widths
% Correct order of tables after \paragraph or \subparagraph
\usepackage{etoolbox}
\makeatletter
\patchcmd\longtable{\par}{\if@noskipsec\mbox{}\fi\par}{}{}
\makeatother
% Allow footnotes in longtable head/foot
\IfFileExists{footnotehyper.sty}{\usepackage{footnotehyper}}{\usepackage{footnote}}
\makesavenoteenv{longtable}
\setlength{\emergencystretch}{3em} % prevent overfull lines
\providecommand{\tightlist}{%
  \setlength{\itemsep}{0pt}\setlength{\parskip}{0pt}}
\setcounter{secnumdepth}{5}
\ifLuaTeX
  \usepackage{selnolig}  % disable illegal ligatures
\fi
\usepackage{bookmark}
\IfFileExists{xurl.sty}{\usepackage{xurl}}{} % add URL line breaks if available
\urlstyle{same}
\hypersetup{
  hidelinks,
  pdfcreator={LaTeX via pandoc}}

\author{}
\date{}

\begin{document}

{
\setcounter{tocdepth}{3}
\tableofcontents
}
\section{V3.3.1 Cost Model 设计说明：Structure-aware Cycle
Correction}\label{v3.3.1-cost-model-ux8bbeux8ba1ux8bf4ux660estructure-aware-cycle-correction}

本文档把当前 cost model 从输入、结构证据、分项
cycles、overlap、sync、penalty 到最终 \texttt{predicted\_cycles}
完整展开。核心目标只有一个：

\begin{Shaded}
\begin{Highlighting}[]
\NormalTok{估计一个候选策略的 total predicted cycles，用于候选排序和可解释分析。}
\end{Highlighting}
\end{Shaded}

当前模型不是实机 msprof cycles，也不是 profiling-calibrated hardware
model。它是一个 \textbf{structure-aware analytical cycle estimation
model}：先用硬件参数和候选策略算出基础 cycles，再用 MLIR
与编译产物中的结构证据修正各分项 cycles 的系统误差，最后汇总为总 cycles
估计。

\begin{center}\rule{0.5\linewidth}{0.5pt}\end{center}

\subsection{1. 一句话理解当前 cost
model}\label{ux4e00ux53e5ux8bddux7406ux89e3ux5f53ux524d-cost-model}

当前版本不要理解成：

\begin{Shaded}
\begin{Highlighting}[]
\NormalTok{score = load + compute + store + sync + 一堆额外加权惩罚}
\end{Highlighting}
\end{Shaded}

而应该理解成：

\begin{Shaded}
\begin{Highlighting}[]
\NormalTok{predicted\_cycles = sum(corrected component cycle estimates) + necessary constraint/risk penalties}
\end{Highlighting}
\end{Shaded}

更直白地说：

\begin{Shaded}
\begin{Highlighting}[]
\NormalTok{第一步：先估计这个策略本来要搬多少数据、算多少 cube/vector、做多少 sync、产生多少 scalar/control 开销。}
\NormalTok{第二步：再看这个 kernel 的 MLIR/产物结构到底偏 memory、compute、vector、scalar 还是 sync。}
\NormalTok{第三步：用这些结构证据修正对应分项的 cycles，而不是到处额外加分扣分。}
\NormalTok{第四步：把修正后的分项 cycles 加起来，得到 total predicted cycles。}
\end{Highlighting}
\end{Shaded}

\begin{center}\rule{0.5\linewidth}{0.5pt}\end{center}

\subsection{2. 候选策略变量}\label{ux5019ux9009ux7b56ux7565ux53d8ux91cf}

一个候选策略记作：

\begin{Shaded}
\begin{Highlighting}[]
\NormalTok{x = (T, M, P, Y)}
\end{Highlighting}
\end{Shaded}

其中：

\begin{longtable}[]{@{}
  >{\raggedright\arraybackslash}p{(\columnwidth - 4\tabcolsep) * \real{0.3333}}
  >{\raggedright\arraybackslash}p{(\columnwidth - 4\tabcolsep) * \real{0.3333}}
  >{\raggedright\arraybackslash}p{(\columnwidth - 4\tabcolsep) * \real{0.3333}}@{}}
\toprule\noalign{}
\begin{minipage}[b]{\linewidth}\raggedright
符号
\end{minipage} & \begin{minipage}[b]{\linewidth}\raggedright
Plan
\end{minipage} & \begin{minipage}[b]{\linewidth}\raggedright
主要控制什么
\end{minipage} \\
\midrule\noalign{}
\endhead
\bottomrule\noalign{}
\endlastfoot
\texttt{T} & \texttt{TilingPlan} &
tile\_m/tile\_n/tile\_k、block\_dim、loop\_order、tail\_strategy、reduce\_split \\
\texttt{M} & \texttt{MultiBufferPlan} & double buffer、per-buffer
multiplier、stage buffer、load/store overlap \\
\texttt{P} & \texttt{CVPipelinePlan} & cube-vector pipeline、pipeline
stage、producer-consumer distance、CV template \\
\texttt{Y} & \texttt{SyncPlan} & keep existing / graph sync
solver、event reuse、sync granularity、sync motion \\
\end{longtable}

cost model 的任务是：给每个 \texttt{x} 算一个
\texttt{predicted\_cycles(x)}，然后搜索器选择 cycles
最低且通过硬件边界检查的候选。

\begin{center}\rule{0.5\linewidth}{0.5pt}\end{center}

\subsection{3. 总公式}\label{ux603bux516cux5f0f}

当前总成本为：

\begin{Shaded}
\begin{Highlighting}[]
\NormalTok{T\_total(x) =}
\NormalTok{    T\_tiles(x)}
\NormalTok{  + T\_sync(x)}
\NormalTok{  + P\_capacity(x)}
\NormalTok{  + P\_shape(x)}
\NormalTok{  + P\_legality(x)}
\end{Highlighting}
\end{Shaded}

其中：

\begin{longtable}[]{@{}
  >{\raggedright\arraybackslash}p{(\columnwidth - 4\tabcolsep) * \real{0.3333}}
  >{\raggedright\arraybackslash}p{(\columnwidth - 4\tabcolsep) * \real{0.3333}}
  >{\raggedright\arraybackslash}p{(\columnwidth - 4\tabcolsep) * \real{0.3333}}@{}}
\toprule\noalign{}
\begin{minipage}[b]{\linewidth}\raggedright
项
\end{minipage} & \begin{minipage}[b]{\linewidth}\raggedright
含义
\end{minipage} & \begin{minipage}[b]{\linewidth}\raggedright
是否是 cycles 语义
\end{minipage} \\
\midrule\noalign{}
\endhead
\bottomrule\noalign{}
\endlastfoot
\texttt{T\_tiles} & 所有 tile 的主计算/搬运/标量控制成本，考虑并行度 &
是，analytical cycles \\
\texttt{T\_sync} & sync/barrier/event 的估计成本 & 是，analytical
cycles \\
\texttt{P\_capacity} & UB/L1/L0/GM workspace 容量压力软惩罚 & 近似
cycles penalty \\
\texttt{P\_shape} & shape/tile 不规整导致的软惩罚 & 近似 cycles
penalty \\
\texttt{P\_legality} & GraphSyncSolver / event reuse / CV pipeline
未验证合法性的风险惩罚 & risk penalty，不等价于真实执行 cycles \\
\end{longtable}

\texttt{T\_tiles} 定义为：

\begin{Shaded}
\begin{Highlighting}[]
\NormalTok{T\_tiles(x) = N\_tiles / E\_parallel * T\_tile(x)}
\end{Highlighting}
\end{Shaded}

其中：

\begin{Shaded}
\begin{Highlighting}[]
\NormalTok{N\_tiles      = 当前 tile shape 下的 tile 数量}
\NormalTok{E\_parallel   = effective\_parallelism}
\NormalTok{T\_tile       = 单 tile 的 corrected steady time}
\end{Highlighting}
\end{Shaded}

\texttt{E\_parallel} 会考虑：

\begin{Shaded}
\begin{Highlighting}[]
\NormalTok{active\_blocks = min(block\_dim, available\_cores, ceil(N\_tiles))}
\NormalTok{waves = ceil(N\_tiles / active\_blocks)}
\NormalTok{tail\_efficiency = N\_tiles / (waves * active\_blocks)}
\NormalTok{E\_parallel = active\_blocks * clamp(tail\_efficiency, 0.20, 1.00)}
\end{Highlighting}
\end{Shaded}

所以 block\_dim 不是越大越好。如果 tile 数很少、tail wave
很碎，有效并行度会被 tail efficiency 拉低。

\begin{center}\rule{0.5\linewidth}{0.5pt}\end{center}

\subsection{4. 单 tile 成本：为什么不是简单 load + compute +
store？}\label{ux5355-tile-ux6210ux672cux4e3aux4ec0ux4e48ux4e0dux662fux7b80ux5355-load-compute-store}

如果没有 double buffer / CV pipeline，模型采用保守串行形式：

\begin{Shaded}
\begin{Highlighting}[]
\NormalTok{T\_tile =}
\NormalTok{    T\_load\_corrected}
\NormalTok{  + T\_compute\_corrected}
\NormalTok{  + T\_store\_corrected}
\NormalTok{  + T\_workspace\_exposed}
\NormalTok{  + T\_scalar\_corrected}
\NormalTok{  + T\_schedule}
\end{Highlighting}
\end{Shaded}

如果开启了 double buffer 或 CV pipeline，load、compute、store
理论上可以部分重叠，所以模型采用 exposed-time 形式：

\begin{Shaded}
\begin{Highlighting}[]
\NormalTok{T\_tile =}
\NormalTok{    max(T\_load\_exposed, T\_compute\_corrected, T\_store\_exposed)}
\NormalTok{  + T\_workspace\_exposed}
\NormalTok{  + T\_scalar\_corrected}
\NormalTok{  + T\_schedule}
\NormalTok{  + T\_warmup\_drain}
\end{Highlighting}
\end{Shaded}

这个设计的含义是：

\begin{Shaded}
\begin{Highlighting}[]
\NormalTok{在稳定流水里，主导时间不是 load + compute + store 全部相加，}
\NormalTok{而是三个 pipeline 中暴露出来的最长路径，加上不能被隐藏的 scalar/control、workspace、schedule 和 warmup/drain。}
\end{Highlighting}
\end{Shaded}

\begin{center}\rule{0.5\linewidth}{0.5pt}\end{center}

\subsection{5. Base
cycles：先算一个没有结构修正的基础估计}\label{base-cyclesux5148ux7b97ux4e00ux4e2aux6ca1ux6709ux7ed3ux6784ux4feeux6b63ux7684ux57faux7840ux4f30ux8ba1}

模型先根据候选策略和硬件参数计算基础项：

\begin{Shaded}
\begin{Highlighting}[]
\NormalTok{T\_load\_base}
\NormalTok{T\_store\_base}
\NormalTok{T\_cube\_base}
\NormalTok{T\_vector\_base}
\NormalTok{T\_fix\_base}
\NormalTok{T\_sync\_base}
\NormalTok{T\_scalar\_base}
\NormalTok{T\_workspace\_base}
\end{Highlighting}
\end{Shaded}

这些基础项来自：

\begin{longtable}[]{@{}
  >{\raggedright\arraybackslash}p{(\columnwidth - 2\tabcolsep) * \real{0.5000}}
  >{\raggedright\arraybackslash}p{(\columnwidth - 2\tabcolsep) * \real{0.5000}}@{}}
\toprule\noalign{}
\begin{minipage}[b]{\linewidth}\raggedright
base 项
\end{minipage} & \begin{minipage}[b]{\linewidth}\raggedright
主要来自
\end{minipage} \\
\midrule\noalign{}
\endhead
\bottomrule\noalign{}
\endlastfoot
\texttt{T\_load\_base} & tile bytes、dtype bytes、GM/UB/L1/L0
读路径、MTE 带宽估计 \\
\texttt{T\_store\_base} & tile 输出 bytes、store/fixpipe path、MTE
写路径估计 \\
\texttt{T\_cube\_base} & tile\_m/tile\_n/tile\_k、cube op proxy、cube
吞吐估计 \\
\texttt{T\_vector\_base} & vector op count proxy、vector tile
workload、vector throughput \\
\texttt{T\_fix\_base} & fixpipe/layout conversion/format conversion
的估计 \\
\texttt{T\_sync\_base} & barrier、set\_flag、wait\_flag、sync\_block
的数量与经验 latency \\
\texttt{T\_scalar\_base} & scalar/control/address/index/cast/loop
调度基础开销 \\
\texttt{T\_workspace\_base} & GM workspace spill/fallback/handoff
traffic \\
\end{longtable}

注意：这些不是实机测量值，而是 analytical
estimate。正因为基础估计会有系统偏差，才需要结构证据修正。

\begin{center}\rule{0.5\linewidth}{0.5pt}\end{center}

\subsection{6.
结构证据到底是什么？}\label{ux7ed3ux6784ux8bc1ux636eux5230ux5e95ux662fux4ec0ux4e48}

当前模型使用三类输入证据：

\begin{Shaded}
\begin{Highlighting}[]
\NormalTok{1. MLIR 静态结构证据}
\NormalTok{2. DES 产物结构证据}
\NormalTok{3. trace event name/count 结构证据}
\end{Highlighting}
\end{Shaded}

它们共同生成 \texttt{kernel\_cost\_profile.json}。

\subsubsection{6.1 MLIR
静态结构证据}\label{mlir-ux9759ux6001ux7ed3ux6784ux8bc1ux636e}

MLIR 侧主要抽取：

\begin{longtable}[]{@{}
  >{\raggedright\arraybackslash}p{(\columnwidth - 4\tabcolsep) * \real{0.3333}}
  >{\raggedright\arraybackslash}p{(\columnwidth - 4\tabcolsep) * \real{0.3333}}
  >{\raggedright\arraybackslash}p{(\columnwidth - 4\tabcolsep) * \real{0.3333}}@{}}
\toprule\noalign{}
\begin{minipage}[b]{\linewidth}\raggedright
证据
\end{minipage} & \begin{minipage}[b]{\linewidth}\raggedright
例子
\end{minipage} & \begin{minipage}[b]{\linewidth}\raggedright
说明
\end{minipage} \\
\midrule\noalign{}
\endhead
\bottomrule\noalign{}
\endlastfoot
flat op counts & \texttt{mmadL1}, \texttt{copy}, \texttt{nd2nz},
\texttt{fixpipe}, vector op, scalar op & 判断 kernel 的原始结构组成 \\
scalar family counts & \texttt{arith\_scalar}, \texttt{index\_cast},
\texttt{pointer\_cast}, \texttt{scf\_for}, \texttt{scf\_if} & 判断
scalar/control/address 开销 \\
sync counts & \texttt{pipe\_barrier}, \texttt{set\_flag},
\texttt{wait\_flag}, \texttt{sync\_block\_set/wait} & 判断同步密度 \\
loop-weighted counts & 内层循环中的 compute/memory/vector/scalar/sync &
内层操作重复执行，权重更高 \\
memory path & GM/UB/L1/L0 之间的静态 path bytes & 判断 memory path
压力 \\
buffer lifetime & buffer live span、byte-span pressure & 判断
workspace/buffer pressure \\
alignment/tail & dim/offset misalignment、mask/tail ops & 判断
vector/fix/scalar fragmentation \\
sequence pattern &
copy-\textgreater nd2nz-\textgreater cube、cube-\textgreater fixpipe-\textgreater vector
& 判断 pipeline 机会 \\
\end{longtable}

\subsubsection{6.2 DES
产物结构证据}\label{des-ux4ea7ux7269ux7ed3ux6784ux8bc1ux636e}

DES 侧主要使用：

\begin{longtable}[]{@{}
  >{\raggedright\arraybackslash}p{(\columnwidth - 2\tabcolsep) * \real{0.5000}}
  >{\raggedright\arraybackslash}p{(\columnwidth - 2\tabcolsep) * \real{0.5000}}@{}}
\toprule\noalign{}
\begin{minipage}[b]{\linewidth}\raggedright
证据
\end{minipage} & \begin{minipage}[b]{\linewidth}\raggedright
使用方式
\end{minipage} \\
\midrule\noalign{}
\endhead
\bottomrule\noalign{}
\endlastfoot
pipe fraction & 判断 lowering 后更像 compute/memory/vector/scalar/sync
哪类 kernel \\
DMA bytes by space path & 增强 memory evidence \\
sync/barrier ops & 增强 sync evidence \\
critical pipe & 作为报告解释字段 \\
\end{longtable}

重要边界：

\begin{Shaded}
\begin{Highlighting}[]
\NormalTok{当前版本不使用 DES makespan/global duration 做 target calibration。}
\end{Highlighting}
\end{Shaded}

也就是说，模型不会做：

\begin{Shaded}
\begin{Highlighting}[]
\NormalTok{predicted\_cycles *= DES\_makespan / base\_prediction}
\end{Highlighting}
\end{Shaded}

它只把 DES pipe mix 当作结构比例证据。例如 pipe\_s 很高，说明 lowering
后 scalar/control 路径很重，于是主要修正 scalar/control cycles。

\subsubsection{6.3 trace event name/count
结构证据}\label{trace-event-namecount-ux7ed3ux6784ux8bc1ux636e}

trace 侧只使用 event name 和 count 的结构提示，例如：

\begin{Shaded}
\begin{Highlighting}[]
\NormalTok{index\_cast/cmpi/addi/muli/load/apply {-}\textgreater{} scalar hint}
\NormalTok{barrier/set\_flag/wait\_flag/sync      {-}\textgreater{} sync hint}
\end{Highlighting}
\end{Shaded}

同样不读取实测 duration target。

\begin{center}\rule{0.5\linewidth}{0.5pt}\end{center}

\subsection{7.
结构比例如何融合？}\label{ux7ed3ux6784ux6bd4ux4f8bux5982ux4f55ux878dux5408}

模型先从 MLIR 得到一组静态比例：

\begin{Shaded}
\begin{Highlighting}[]
\NormalTok{static\_ratios = normalize(\{compute\_score, memory\_score, vector\_score, scalar\_score, sync\_score\})}
\end{Highlighting}
\end{Shaded}

如果 DES 产物可用，再得到一组产物比例：

\begin{Shaded}
\begin{Highlighting}[]
\NormalTok{product\_ratios = normalize(\{compute\_pipe, memory\_pipe, vector\_pipe, scalar\_pipe, sync\_pipe\})}
\end{Highlighting}
\end{Shaded}

当前 V3.3.1 的融合方式为：

\begin{Shaded}
\begin{Highlighting}[]
\NormalTok{structure\_ratios = 0.60 * static\_ratios + 0.40 * product\_ratios}
\NormalTok{structure\_ratios = normalize(structure\_ratios)}
\end{Highlighting}
\end{Shaded}

为什么不是产物 65\% 或 80\%？因为 DES pipe fraction 不是实机
profiling。它能反映 lowering 后结构，但仍可能被某一个 pipe
字段支配。降低产物权重可以避免单样本产物过度主导 cost model。

在你上传的样本中，融合后 profile 类似：

\begin{Shaded}
\begin{Highlighting}[]
\NormalTok{compute = 0.016}
\NormalTok{memory  = 0.116}
\NormalTok{vector  = 0.079}
\NormalTok{scalar  = 0.553}
\NormalTok{sync    = 0.236}
\end{Highlighting}
\end{Shaded}

因此该 kernel 被判为：

\begin{Shaded}
\begin{Highlighting}[]
\NormalTok{kernel\_type = scalar\_control\_heavy}
\NormalTok{dominant\_component = scalar}
\end{Highlighting}
\end{Shaded}

这不是说真实 cycles 中 scalar 一定占 55.3\%，而是说结构证据显示
scalar/control 是最强的修正方向。

\begin{center}\rule{0.5\linewidth}{0.5pt}\end{center}

\subsection{8. 从结构比例到 cycle correction
factors}\label{ux4eceux7ed3ux6784ux6bd4ux4f8bux5230-cycle-correction-factors}

V3.3.1 的核心改动是：结构证据不再到处扩散成一堆
reward/penalty，而是映射为分项 cycle correction factors。

当前主要 factors 为：

\begin{Shaded}
\begin{Highlighting}[]
\NormalTok{m\_mem       = memory\_cycle\_correction}
\NormalTok{m\_mem\_path  = memory\_path\_cycle\_correction}
\NormalTok{m\_cube      = compute\_cycle\_correction}
\NormalTok{m\_vec       = vector\_cycle\_correction}
\NormalTok{m\_align     = alignment\_cycle\_correction}
\NormalTok{m\_scalar    = scalar\_cycle\_correction}
\NormalTok{m\_frag      = small\_tile\_fragmentation\_correction}
\NormalTok{m\_sync      = sync\_cycle\_correction}
\NormalTok{m\_ws        = workspace\_pressure\_correction}
\NormalTok{c\_overlap   = overlap\_confidence}
\NormalTok{c\_cv         = cv\_overlap\_confidence}
\end{Highlighting}
\end{Shaded}

\subsubsection{8.1 每个 factor
的语义}\label{ux6bcfux4e2a-factor-ux7684ux8bedux4e49}

为了避免 PDF 表格过宽，下面逐项说明每个 factor 的含义：

\begin{itemize}
\tightlist
\item
  \texttt{memory\_cycle\_correction}：修正 load/store/fix 的 memory
  side。主要来自 \texttt{memory\_ratio}、memory path bytes、buffer
  pressure。含义是搬运路径比基础估计更重或更轻。
\item
  \texttt{memory\_path\_cycle\_correction}：修正 load/store。主要来自
  GM/UB/L1/L0 path bytes。含义是具体 memory path 的复杂度修正。
\item
  \texttt{compute\_cycle\_correction}：修正 cube compute。主要来自
  compute ratio 和 CV opportunity。它不是 cube reward，只是 cube cycles
  的估计修正。
\item
  \texttt{vector\_cycle\_correction}：修正 vector compute。主要来自
  vector ratio、alignment/tail proxy。
\item
  \texttt{alignment\_cycle\_correction}：修正 vector/fix/scalar
  fragmentation。主要来自 dim/offset misalignment 和 mask/tail ops。
\item
  \texttt{scalar\_cycle\_correction}：修正 scalar/control
  cycles。主要来自 scalar ratio、loop-weighted scalar ratio、alignment
  proxy。
\item
  \texttt{small\_tile\_fragmentation\_correction}：修正小 tile 下的
  scalar/control 碎片化成本。主要来自 tile fragmentation 和 loop
  scalar。
\item
  \texttt{sync\_cycle\_correction}：修正 sync cost。主要来自 sync ratio
  和 sync criticality。
\item
  \texttt{workspace\_pressure\_correction}：修正 workspace exposed
  cycles。主要来自 buffer lifetime pressure。
\item
  \texttt{overlap\_confidence}：修正 load/store overlap
  的可信度。主要来自 scalar/sync/memory/CV
  opportunity，但只做窄范围修正。
\item
  \texttt{cv\_overlap\_confidence}：修正 cube-vector overlap
  的可信度。主要来自 compute/vector balance、scalar/sync
  density，也只做窄范围修正。
\end{itemize}

\subsubsection{8.2 设计边界}\label{ux8bbeux8ba1ux8fb9ux754c}

\begin{Shaded}
\begin{Highlighting}[]
\NormalTok{memory evidence 主要进 memory cycles}
\NormalTok{scalar evidence 主要进 scalar/control cycles}
\NormalTok{sync evidence 主要进 sync cycles}
\NormalTok{alignment evidence 主要进 vector/fix/scalar fragmentation}
\NormalTok{overlap confidence 只做窄范围修正}
\NormalTok{legality risk 单独进入 P\_legality}
\end{Highlighting}
\end{Shaded}

这个边界是为了避免 double counting。

\begin{center}\rule{0.5\linewidth}{0.5pt}\end{center}

\subsection{9. Corrected component cycles
逐项展开}\label{corrected-component-cycles-ux9010ux9879ux5c55ux5f00}

\subsubsection{9.1 Load / Store}\label{load-store}

\begin{Shaded}
\begin{Highlighting}[]
\NormalTok{T\_load\_corrected  = T\_load\_base  * m\_mem * m\_mem\_path}
\NormalTok{T\_store\_corrected = T\_store\_base * m\_mem * m\_mem\_path}
\end{Highlighting}
\end{Shaded}

解释：

\begin{Shaded}
\begin{Highlighting}[]
\NormalTok{基础 load/store 根据 bytes / bandwidth 估计；}
\NormalTok{如果 MLIR/产物显示 memory path 更复杂、DMA path 更多、buffer pressure 更高，}
\NormalTok{则通过 m\_mem 和 m\_mem\_path 修正搬运 cycles。}
\end{Highlighting}
\end{Shaded}

\subsubsection{9.2 Cube / Vector / Fix}\label{cube-vector-fix}

\begin{Shaded}
\begin{Highlighting}[]
\NormalTok{T\_cube\_corrected   = T\_cube\_base   * m\_cube}
\NormalTok{T\_vector\_corrected = T\_vector\_base * m\_vec * m\_align}
\NormalTok{T\_fix\_corrected    = T\_fix\_base    * m\_mem * m\_align}
\end{Highlighting}
\end{Shaded}

解释：

\begin{Shaded}
\begin{Highlighting}[]
\NormalTok{cube 主要受 compute correction 修正；}
\NormalTok{vector 受 vector correction 和 alignment correction 修正；}
\NormalTok{fixpipe/layout conversion 同时带有 memory side 和 alignment side，因此使用 m\_mem 与 m\_align。}
\end{Highlighting}
\end{Shaded}

\subsubsection{9.3 Cube-vector pipeline}\label{cube-vector-pipeline}

\begin{Shaded}
\begin{Highlighting}[]
\NormalTok{T\_compute\_corrected =}
\NormalTok{    T\_cube\_corrected}
\NormalTok{  + T\_vector\_corrected}
\NormalTok{  {-} r\_cv * c\_cv * min(T\_cube\_corrected, T\_vector\_corrected)}
\NormalTok{  + T\_fix\_corrected}
\end{Highlighting}
\end{Shaded}

其中：

\begin{Shaded}
\begin{Highlighting}[]
\NormalTok{r\_cv = CVPipelinePlan 给出的策略 overlap ratio}
\NormalTok{c\_cv = 结构证据给出的 CV overlap confidence}
\end{Highlighting}
\end{Shaded}

解释：

\begin{Shaded}
\begin{Highlighting}[]
\NormalTok{CV pipeline 的主要收益来自候选策略本身；}
\NormalTok{结构证据只判断这个收益是否可信，不能把 compute{-}heavy 直接变成额外 reward。}
\end{Highlighting}
\end{Shaded}

\subsubsection{9.4 Scalar/control}\label{scalarcontrol}

\begin{Shaded}
\begin{Highlighting}[]
\NormalTok{T\_scalar\_corrected = T\_scalar\_base * m\_scalar * m\_frag}
\end{Highlighting}
\end{Shaded}

解释：

\begin{Shaded}
\begin{Highlighting}[]
\NormalTok{scalar/control 不是主 cube/vector pipeline 的一部分；}
\NormalTok{大量 index\_cast、pointer\_cast、scf\_for/scf\_if、get\_block\_idx、mask/tail、inner{-}loop scalar op 会形成不可忽视的调度/控制成本。}
\end{Highlighting}
\end{Shaded}

在你上传的样本里，scalar/control 证据很强，所以
\texttt{scalar\_cycle\_correction} 明显高于 1。这是合理的，因为该 kernel
的 MLIR 和产物都显示 scalar/control-heavy。

\subsubsection{9.5 Workspace}\label{workspace}

\begin{Shaded}
\begin{Highlighting}[]
\NormalTok{T\_workspace\_exposed = T\_workspace\_base * m\_ws}
\end{Highlighting}
\end{Shaded}

解释：

\begin{Shaded}
\begin{Highlighting}[]
\NormalTok{GM workspace / handoff / spill traffic 不是普通 load/store 的完全可隐藏部分；}
\NormalTok{它更像 fallback traffic，因此作为 exposed cost 叠加。}
\end{Highlighting}
\end{Shaded}

\subsubsection{9.6 Sync}\label{sync}

\begin{Shaded}
\begin{Highlighting}[]
\NormalTok{T\_sync = T\_sync\_base * m\_sync}
\end{Highlighting}
\end{Shaded}

其中基础 sync 通常来自：

\begin{Shaded}
\begin{Highlighting}[]
\NormalTok{T\_sync\_base =}
\NormalTok{    num\_barrier * barrier\_latency}
\NormalTok{  + num\_set\_flag * set\_flag\_latency}
\NormalTok{  + num\_wait\_flag * wait\_flag\_latency}
\NormalTok{  + sync\_block\_cost}
\end{Highlighting}
\end{Shaded}

解释：

\begin{Shaded}
\begin{Highlighting}[]
\NormalTok{sync evidence 只修正同步操作本身的估计开销。}
\NormalTok{GraphSyncSolver 未验证是否合法，不应该混进 T\_sync，而应该放到 P\_legality。}
\end{Highlighting}
\end{Shaded}

\begin{center}\rule{0.5\linewidth}{0.5pt}\end{center}

\subsection{10. Overlap
如何处理？}\label{overlap-ux5982ux4f55ux5904ux7406}

\subsubsection{10.1 Load/store overlap}\label{loadstore-overlap}

候选策略先给出基础 overlap：

\begin{Shaded}
\begin{Highlighting}[]
\NormalTok{r\_load\_strategy}
\NormalTok{r\_store\_strategy}
\end{Highlighting}
\end{Shaded}

结构证据只给出 confidence：

\begin{Shaded}
\begin{Highlighting}[]
\NormalTok{r\_load  = r\_load\_strategy  * c\_overlap}
\NormalTok{r\_store = r\_store\_strategy * c\_overlap}
\end{Highlighting}
\end{Shaded}

然后：

\begin{Shaded}
\begin{Highlighting}[]
\NormalTok{T\_load\_exposed  = T\_load\_corrected  * (1 {-} r\_load)}
\NormalTok{T\_store\_exposed = T\_store\_corrected * (1 {-} r\_store)}
\end{Highlighting}
\end{Shaded}

\subsubsection{10.2 为什么 overlap confidence
是窄范围？}\label{ux4e3aux4ec0ux4e48-overlap-confidence-ux662fux7a84ux8303ux56f4}

因为 overlap 的主因应该是策略和硬件合法性，例如：

\begin{Shaded}
\begin{Highlighting}[]
\NormalTok{double\_buffer 是否开启}
\NormalTok{per{-}buffer multiplier 是否足够}
\NormalTok{stage buffer 是否放得下}
\NormalTok{producer{-}consumer distance 是否合理}
\NormalTok{UB/L1/L0 是否 overflow}
\end{Highlighting}
\end{Shaded}

结构证据只能说``这个 overlap 可信度更高或更低''，不能说``scalar-heavy
所有 overlap 全部大幅失效''。因此 V3.3.1 将 overlap confidence
限定在较窄范围，避免重复惩罚。

\begin{center}\rule{0.5\linewidth}{0.5pt}\end{center}

\subsection{11. Penalty 的语义}\label{penalty-ux7684ux8bedux4e49}

\subsubsection{11.1 Capacity penalty}\label{capacity-penalty}

\begin{Shaded}
\begin{Highlighting}[]
\NormalTok{P\_capacity = memory\_pressure\_penalty(scope\_utilization, hw)}
\end{Highlighting}
\end{Shaded}

它表示：虽然没有触发硬件 gate 的 hard overflow，但 UB/L1/L0
等资源已经接近边界，可能导致更高风险或更差调度。

\subsubsection{11.2 Shape penalty}\label{shape-penalty}

\begin{Shaded}
\begin{Highlighting}[]
\NormalTok{P\_shape = shape\_regularization\_penalty(tile shape, kernel shape, hw)}
\end{Highlighting}
\end{Shaded}

它表示：tile shape 与 kernel shape/hardware granularity 不匹配，可能产生
tail、fragmentation 或不规则访问。

\subsubsection{11.3 Legality risk penalty}\label{legality-risk-penalty}

\begin{Shaded}
\begin{Highlighting}[]
\NormalTok{P\_legality =}
\NormalTok{    P\_graph\_sync\_unknown}
\NormalTok{  + P\_event\_reuse\_unknown}
\NormalTok{  + P\_cv\_pipeline\_estimated}
\end{Highlighting}
\end{Shaded}

它专门表达策略合法性/可落地性不确定性。例如：

\begin{Shaded}
\begin{Highlighting}[]
\NormalTok{GraphSyncSolver status = UNKNOWN}
\NormalTok{CV pipeline legality = PASS\_ESTIMATED}
\end{Highlighting}
\end{Shaded}

这类风险不能说是真实执行时间，因此单独放在 \texttt{P\_legality}，不污染
\texttt{T\_sync} 的 cycles 语义。

\begin{center}\rule{0.5\linewidth}{0.5pt}\end{center}

\subsection{12. V3.3.1
如何利用结构化信息？}\label{v3.3.1-ux5982ux4f55ux5229ux7528ux7ed3ux6784ux5316ux4fe1ux606f}

完整链路如下：

\begin{Shaded}
\begin{Highlighting}[]
\NormalTok{MLIR + product artifacts}
\NormalTok{        |}
\NormalTok{        v}
\NormalTok{extract structural evidence}
\NormalTok{        |}
\NormalTok{        v}
\NormalTok{compute static\_ratios and product\_ratios}
\NormalTok{        |}
\NormalTok{        v}
\NormalTok{fuse ratios: 60\% MLIR + 40\% artifact}
\NormalTok{        |}
\NormalTok{        v}
\NormalTok{build kernel\_cost\_profile}
\NormalTok{        |}
\NormalTok{        v}
\NormalTok{map evidence to cycle correction factors}
\NormalTok{        |}
\NormalTok{        v}
\NormalTok{correct load/compute/store/scalar/sync/workspace cycles}
\NormalTok{        |}
\NormalTok{        v}
\NormalTok{apply strategy overlap and narrow confidence correction}
\NormalTok{        |}
\NormalTok{        v}
\NormalTok{sum corrected cycles + necessary penalties}
\NormalTok{        |}
\NormalTok{        v}
\NormalTok{predicted\_cycles}
\end{Highlighting}
\end{Shaded}

也就是：结构化信息现在不是独立 score，而是进入每个分项 cycle estimate。

\subsubsection{12.1 例子：一个 scalar-heavy
kernel}\label{ux4f8bux5b50ux4e00ux4e2a-scalar-heavy-kernel}

如果 MLIR/产物显示：

\begin{Shaded}
\begin{Highlighting}[]
\NormalTok{scalar\_ratio 高}
\NormalTok{loop\_weighted\_scalar 高}
\NormalTok{pipe\_s fraction 高}
\NormalTok{sync density 较高}
\end{Highlighting}
\end{Shaded}

V3.3.1 的行为是：

\begin{Shaded}
\begin{Highlighting}[]
\NormalTok{scalar\_cycle\_correction 上升}
\NormalTok{small\_tile\_fragmentation\_correction 上升}
\NormalTok{sync\_cycle\_correction 适度上升}
\NormalTok{overlap\_confidence 轻微下降}
\end{Highlighting}
\end{Shaded}

但它不会再让同一个 scalar-heavy 证据同时大幅：

\begin{Shaded}
\begin{Highlighting}[]
\NormalTok{提高 scalar cost}
\NormalTok{提高 sync cost}
\NormalTok{提高 small tile penalty}
\NormalTok{大幅压低 overlap reward}
\NormalTok{大幅压低 CV reward}
\NormalTok{提高 legality risk}
\end{Highlighting}
\end{Shaded}

这就是 V3.3.1 相比 V3.3 更干净的地方。

\begin{center}\rule{0.5\linewidth}{0.5pt}\end{center}

\subsection{13. 和 V3.3 的区别}\label{ux548c-v3.3-ux7684ux533aux522b}

\subsubsection{13.1 V3.3 的逻辑}\label{v3.3-ux7684ux903bux8f91}

V3.3 更像：

\begin{Shaded}
\begin{Highlighting}[]
\NormalTok{结构证据 {-}\textgreater{} 多个 multiplier/reward/penalty {-}\textgreater{} total predicted cycles}
\end{Highlighting}
\end{Shaded}

当 kernel 被判为 scalar/sync-heavy 时，同一类证据可能同时影响：

\begin{Shaded}
\begin{Highlighting}[]
\NormalTok{scalar\_control\_multiplier}
\NormalTok{small\_tile\_scalar\_penalty\_scale}
\NormalTok{loop\_weighted\_scalar\_multiplier}
\NormalTok{sync\_multiplier}
\NormalTok{sync\_criticality\_multiplier}
\NormalTok{overlap\_reward\_scale}
\NormalTok{cv\_reward\_scale}
\NormalTok{cube\_reward\_scale}
\NormalTok{legality\_risk\_penalty}
\end{Highlighting}
\end{Shaded}

问题是：这些路径之间存在语义重叠。比如 scalar-heavy 既提高
scalar\_control\_time，又降低 overlap，又降低 CV reward，还可能提高
sync/risk。最终 \texttt{predicted\_cycles} 仍然叫
cycles，但内部更像``多重加权 ranking score''。

\subsubsection{13.2 V3.3.1 的逻辑}\label{v3.3.1-ux7684ux903bux8f91}

V3.3.1 改成：

\begin{Shaded}
\begin{Highlighting}[]
\NormalTok{结构证据 {-}\textgreater{} 对应分项 cycle correction {-}\textgreater{} total predicted cycles}
\end{Highlighting}
\end{Shaded}

核心变化：

\begin{longtable}[]{@{}
  >{\raggedright\arraybackslash}p{(\columnwidth - 4\tabcolsep) * \real{0.3333}}
  >{\raggedright\arraybackslash}p{(\columnwidth - 4\tabcolsep) * \real{0.3333}}
  >{\raggedright\arraybackslash}p{(\columnwidth - 4\tabcolsep) * \real{0.3333}}@{}}
\toprule\noalign{}
\begin{minipage}[b]{\linewidth}\raggedright
方面
\end{minipage} & \begin{minipage}[b]{\linewidth}\raggedright
V3.3
\end{minipage} & \begin{minipage}[b]{\linewidth}\raggedright
V3.3.1
\end{minipage} \\
\midrule\noalign{}
\endhead
\bottomrule\noalign{}
\endlastfoot
总体语义 & structure-aware weighted cost & structure-aware corrected
cycle estimate \\
结构证据用途 & 可同时影响多个 cost/reward/penalty &
每类证据主要修正对应分项 cycles \\
scalar evidence & scalar、small tile、sync、overlap、CV reward
多路径影响 & 主要进入 scalar/control 和 fragmentation \\
sync evidence & sync multiplier、criticality multiplier、risk 可能叠加 &
sync cycles 与 legality risk 分离 \\
overlap & 可能被 scalar/sync 较强打折 & 只做窄范围 confidence 修正 \\
DES artifact 权重 & 产物比例更强，容易主导 & MLIR 60\% + artifact 40\%
保守融合 \\
predicted\_cycles 语义 & 更像 cycles-shaped ranking score & 更像
corrected component cycles sum \\
double counting 风险 & 较高 & 明显降低 \\
\end{longtable}

\subsubsection{13.3 为什么 V3.3.1
更适合汇报？}\label{ux4e3aux4ec0ux4e48-v3.3.1-ux66f4ux9002ux5408ux6c47ux62a5}

因为可以清楚回答领导的几个问题：

\begin{itemize}
\tightlist
\item
  问：你的目标到底是算 score 还是 cycles？ 答：目标是估计 total
  predicted cycles。
\item
  问：结构信息干什么用？ 答：修正
  load/compute/store/scalar/sync/workspace 各分项 cycles
  的基础估计误差。
\item
  问：有没有用实机数据硬拟合？ 答：没有。在线寻优不使用 msprof
  target，不使用 DES makespan/global scale。
\item
  问：为什么不是 load + compute + store + sync 直接加？ 答：因为有
  pipeline overlap。主 tile 部分在流水模式下使用
  \texttt{max(exposed\ load,\ compute,\ exposed\ store)}，然后再加
  scalar/control、workspace、schedule 等不可隐藏项。
\item
  问：GraphSyncSolver 不确定怎么办？
  答：执行时间估计和合法性风险分开；未知合法性进入
  \texttt{P\_legality}，不污染 sync cycles。
\end{itemize}

\begin{center}\rule{0.5\linewidth}{0.5pt}\end{center}

\subsection{14.
当前样本的解释方式}\label{ux5f53ux524dux6837ux672cux7684ux89e3ux91caux65b9ux5f0f}

以
\texttt{kernel\_001.npuir.mlir\ +\ prefill\_des.json\ +\ prefill\_trace.json}
为例，当前 profile 显示：

\begin{Shaded}
\begin{Highlighting}[]
\NormalTok{dominant\_component = scalar}
\NormalTok{kernel\_type = scalar\_control\_heavy}
\end{Highlighting}
\end{Shaded}

融合后的结构比例大致为：

\begin{Shaded}
\begin{Highlighting}[]
\NormalTok{compute: 0.016}
\NormalTok{memory : 0.116}
\NormalTok{vector : 0.079}
\NormalTok{scalar : 0.553}
\NormalTok{sync   : 0.236}
\end{Highlighting}
\end{Shaded}

对应 correction factors 大致为：

\begin{Shaded}
\begin{Highlighting}[]
\NormalTok{memory\_cycle\_correction = 1.10}
\NormalTok{compute\_cycle\_correction = 1.06}
\NormalTok{vector\_cycle\_correction = 1.02}
\NormalTok{scalar\_cycle\_correction = 1.72}
\NormalTok{sync\_cycle\_correction = 1.36}
\NormalTok{overlap\_confidence = 0.88}
\NormalTok{cv\_overlap\_confidence = 0.90}
\end{Highlighting}
\end{Shaded}

可以这样解释：

\begin{Shaded}
\begin{Highlighting}[]
\NormalTok{这个 kernel 不是典型 cube{-}heavy kernel，而是 scalar/control 与 sync 结构比较重。}
\NormalTok{因此模型主要提高 scalar/control cycles 和 sync cycles 的估计；}
\NormalTok{memory/compute/vector 只做轻微修正；}
\NormalTok{overlap 与 CV pipeline 收益只被轻度降置信，而不是被结构证据大幅砍掉。}
\end{Highlighting}
\end{Shaded}

这比 V3.3 的解释更稳，因为不会出现``一个 scalar-heavy
判断到处乘''的问题。

\begin{center}\rule{0.5\linewidth}{0.5pt}\end{center}

\subsection{15.
当前模型仍然不能证明什么？}\label{ux5f53ux524dux6a21ux578bux4ecdux7136ux4e0dux80fdux8bc1ux660eux4ec0ux4e48}

当前模型可以用于：

\begin{Shaded}
\begin{Highlighting}[]
\NormalTok{候选策略排序}
\NormalTok{缺陷 MLIR 的相对诊断}
\NormalTok{硬件边界检查}
\NormalTok{解释为什么某些策略被惩罚或被选中}
\NormalTok{在没有实机 profiling 时做 conservative analytical search}
\end{Highlighting}
\end{Shaded}

但不能直接证明：

\begin{Shaded}
\begin{Highlighting}[]
\NormalTok{真实 msprof cycles 一定等于 predicted\_cycles}
\NormalTok{真实 speedup 一定等于 predicted speedup}
\NormalTok{GraphSyncSolver 策略一定可 rewrite 且 deadlock{-}free}
\NormalTok{DES pipe fraction 与实机 pipe active/stall cycles 完全一致}
\end{Highlighting}
\end{Shaded}

如果要把它升级为更真实的 hardware performance model，需要：

\begin{Shaded}
\begin{Highlighting}[]
\NormalTok{同一批 kernel 的 current IR 与候选策略 rewrite 后实机 msprof 数据}
\NormalTok{每个候选的编译产物 / DES 产物，而不是共享一个 kernel{-}level artifact profile}
\NormalTok{用实机数据离线拟合 correction factors 的 alpha/clamp/config}
\NormalTok{用 held{-}out kernel 验证 ranking accuracy 和 cycles prediction error}
\end{Highlighting}
\end{Shaded}

\begin{center}\rule{0.5\linewidth}{0.5pt}\end{center}

\subsection{16.
推荐的口头表述}\label{ux63a8ux8350ux7684ux53e3ux5934ux8868ux8ff0}

汇报时可以这样说：

我们当前的 cost model 目标仍然是估计总 cycles，但不是直接使用实机
profiling。

模型先根据 tile、buffer、pipeline、sync plan 和硬件配置，计算
load、compute、store、scalar/control、sync 等基础 cycles。

然后利用 MLIR 和编译产物中的结构证据生成 kernel-level cycle correction
factors。这些结构证据包括 op count、loop-weighted scalar/sync、memory
path、buffer lifetime、DES pipe mix 和 trace event name count。

这些 correction factors 只修正对应的分项 cycles。例如 memory 证据修正
load/store，scalar 证据修正 scalar/control，sync 证据修正 sync。

最后模型汇总 corrected component cycles，并额外加入容量、shape
和合法性风险惩罚，得到 predicted\_cycles。

相比 V3.3，V3.3.1
最大变化是降低了结构证据的重复加权风险。结构信息不再同时作为多路
reward/penalty 扩散，而是回到分项 cycles 修正的语义上。

\begin{center}\rule{0.5\linewidth}{0.5pt}\end{center}

\subsection{17.
文件和报告中应该怎么看}\label{ux6587ux4ef6ux548cux62a5ux544aux4e2dux5e94ux8be5ux600eux4e48ux770b}

运行后重点看这些文件：

\begin{longtable}[]{@{}
  >{\raggedright\arraybackslash}p{(\columnwidth - 2\tabcolsep) * \real{0.5000}}
  >{\raggedright\arraybackslash}p{(\columnwidth - 2\tabcolsep) * \real{0.5000}}@{}}
\toprule\noalign{}
\begin{minipage}[b]{\linewidth}\raggedright
文件
\end{minipage} & \begin{minipage}[b]{\linewidth}\raggedright
看什么
\end{minipage} \\
\midrule\noalign{}
\endhead
\bottomrule\noalign{}
\endlastfoot
\texttt{kernel\_cost\_profile.json} & kernel 类型、结构比例、cycle
correction factors \\
\texttt{cost\_breakdown.json} & selected strategy 的分项 cycles、overlap
saving、penalty \\
\texttt{selected\_strategy.json} & 最终选择的候选参数 \\
\texttt{selected\_plan.json} & 四类 plan 的具体取值与 legality 状态 \\
\texttt{top\_candidates.json} & Top 候选的排序与差异 \\
\texttt{hardware\_boundary\_audit.json} & UB/L1/L0/GM workspace
是否接近或超过边界 \\
\end{longtable}

最关键的是：不要只看 \texttt{predicted\_cycles} 一个数字，要同时看：

\begin{Shaded}
\begin{Highlighting}[]
\NormalTok{load\_exposed}
\NormalTok{store\_exposed}
\NormalTok{cube\_vector\_time}
\NormalTok{scalar\_control\_time}
\NormalTok{sync\_cost}
\NormalTok{memory\_pressure\_penalty}
\NormalTok{shape\_regularization\_penalty}
\NormalTok{legality\_risk\_penalty}
\NormalTok{kernel\_cost\_profile\_weights}
\end{Highlighting}
\end{Shaded}

这些分项能说明为什么候选被选中或被压下去。

\begin{center}\rule{0.5\linewidth}{0.5pt}\end{center}

\subsection{18. 小结}\label{ux5c0fux7ed3}

V3.3.1 的 cost model 可以概括为：

\begin{Shaded}
\begin{Highlighting}[]
\NormalTok{structure{-}aware corrected component cycle model}
\end{Highlighting}
\end{Shaded}

它的关键优点是：

\begin{Shaded}
\begin{Highlighting}[]
\NormalTok{1. 目标仍然是 total predicted cycles；}
\NormalTok{2. 结构证据用于修正分项 cycles，而不是额外打分；}
\NormalTok{3. memory/scalar/sync/vector/compute 各有相对清晰的修正边界；}
\NormalTok{4. overlap 由策略主导，结构证据只做窄范围 confidence 修正；}
\NormalTok{5. legality risk 与 sync cycles 分离；}
\NormalTok{6. 相比 V3.3，double counting 风险更低，汇报语义更清楚。}
\end{Highlighting}
\end{Shaded}

一句话：

V3.3.1 不是把结构信息拿来``再加权一次总成本''。它的做法是：

\begin{Shaded}
\begin{Highlighting}[]
\NormalTok{结构信息 {-}\textgreater{} 修正每个分项 cycles 的估计误差 {-}\textgreater{} 汇总为 predicted total cycles}
\end{Highlighting}
\end{Shaded}


\end{document}
